% Options for packages loaded elsewhere
\PassOptionsToPackage{unicode}{hyperref}
\PassOptionsToPackage{hyphens}{url}
\documentclass[
]{article}
\usepackage{xcolor}
\usepackage[margin=1in]{geometry}
\usepackage{amsmath,amssymb}
\setcounter{secnumdepth}{-\maxdimen} % remove section numbering
\usepackage{iftex}
\ifPDFTeX
  \usepackage[T1]{fontenc}
  \usepackage[utf8]{inputenc}
  \usepackage{textcomp} % provide euro and other symbols
\else % if luatex or xetex
  \usepackage{unicode-math} % this also loads fontspec
  \defaultfontfeatures{Scale=MatchLowercase}
  \defaultfontfeatures[\rmfamily]{Ligatures=TeX,Scale=1}
\fi
\usepackage{lmodern}
\ifPDFTeX\else
  % xetex/luatex font selection
\fi
% Use upquote if available, for straight quotes in verbatim environments
\IfFileExists{upquote.sty}{\usepackage{upquote}}{}
\IfFileExists{microtype.sty}{% use microtype if available
  \usepackage[]{microtype}
  \UseMicrotypeSet[protrusion]{basicmath} % disable protrusion for tt fonts
}{}
\makeatletter
\@ifundefined{KOMAClassName}{% if non-KOMA class
  \IfFileExists{parskip.sty}{%
    \usepackage{parskip}
  }{% else
    \setlength{\parindent}{0pt}
    \setlength{\parskip}{6pt plus 2pt minus 1pt}}
}{% if KOMA class
  \KOMAoptions{parskip=half}}
\makeatother
\usepackage{color}
\usepackage{fancyvrb}
\newcommand{\VerbBar}{|}
\newcommand{\VERB}{\Verb[commandchars=\\\{\}]}
\DefineVerbatimEnvironment{Highlighting}{Verbatim}{commandchars=\\\{\}}
% Add ',fontsize=\small' for more characters per line
\usepackage{framed}
\definecolor{shadecolor}{RGB}{248,248,248}
\newenvironment{Shaded}{\begin{snugshade}}{\end{snugshade}}
\newcommand{\AlertTok}[1]{\textcolor[rgb]{0.94,0.16,0.16}{#1}}
\newcommand{\AnnotationTok}[1]{\textcolor[rgb]{0.56,0.35,0.01}{\textbf{\textit{#1}}}}
\newcommand{\AttributeTok}[1]{\textcolor[rgb]{0.13,0.29,0.53}{#1}}
\newcommand{\BaseNTok}[1]{\textcolor[rgb]{0.00,0.00,0.81}{#1}}
\newcommand{\BuiltInTok}[1]{#1}
\newcommand{\CharTok}[1]{\textcolor[rgb]{0.31,0.60,0.02}{#1}}
\newcommand{\CommentTok}[1]{\textcolor[rgb]{0.56,0.35,0.01}{\textit{#1}}}
\newcommand{\CommentVarTok}[1]{\textcolor[rgb]{0.56,0.35,0.01}{\textbf{\textit{#1}}}}
\newcommand{\ConstantTok}[1]{\textcolor[rgb]{0.56,0.35,0.01}{#1}}
\newcommand{\ControlFlowTok}[1]{\textcolor[rgb]{0.13,0.29,0.53}{\textbf{#1}}}
\newcommand{\DataTypeTok}[1]{\textcolor[rgb]{0.13,0.29,0.53}{#1}}
\newcommand{\DecValTok}[1]{\textcolor[rgb]{0.00,0.00,0.81}{#1}}
\newcommand{\DocumentationTok}[1]{\textcolor[rgb]{0.56,0.35,0.01}{\textbf{\textit{#1}}}}
\newcommand{\ErrorTok}[1]{\textcolor[rgb]{0.64,0.00,0.00}{\textbf{#1}}}
\newcommand{\ExtensionTok}[1]{#1}
\newcommand{\FloatTok}[1]{\textcolor[rgb]{0.00,0.00,0.81}{#1}}
\newcommand{\FunctionTok}[1]{\textcolor[rgb]{0.13,0.29,0.53}{\textbf{#1}}}
\newcommand{\ImportTok}[1]{#1}
\newcommand{\InformationTok}[1]{\textcolor[rgb]{0.56,0.35,0.01}{\textbf{\textit{#1}}}}
\newcommand{\KeywordTok}[1]{\textcolor[rgb]{0.13,0.29,0.53}{\textbf{#1}}}
\newcommand{\NormalTok}[1]{#1}
\newcommand{\OperatorTok}[1]{\textcolor[rgb]{0.81,0.36,0.00}{\textbf{#1}}}
\newcommand{\OtherTok}[1]{\textcolor[rgb]{0.56,0.35,0.01}{#1}}
\newcommand{\PreprocessorTok}[1]{\textcolor[rgb]{0.56,0.35,0.01}{\textit{#1}}}
\newcommand{\RegionMarkerTok}[1]{#1}
\newcommand{\SpecialCharTok}[1]{\textcolor[rgb]{0.81,0.36,0.00}{\textbf{#1}}}
\newcommand{\SpecialStringTok}[1]{\textcolor[rgb]{0.31,0.60,0.02}{#1}}
\newcommand{\StringTok}[1]{\textcolor[rgb]{0.31,0.60,0.02}{#1}}
\newcommand{\VariableTok}[1]{\textcolor[rgb]{0.00,0.00,0.00}{#1}}
\newcommand{\VerbatimStringTok}[1]{\textcolor[rgb]{0.31,0.60,0.02}{#1}}
\newcommand{\WarningTok}[1]{\textcolor[rgb]{0.56,0.35,0.01}{\textbf{\textit{#1}}}}
\usepackage{graphicx}
\makeatletter
\newsavebox\pandoc@box
\newcommand*\pandocbounded[1]{% scales image to fit in text height/width
  \sbox\pandoc@box{#1}%
  \Gscale@div\@tempa{\textheight}{\dimexpr\ht\pandoc@box+\dp\pandoc@box\relax}%
  \Gscale@div\@tempb{\linewidth}{\wd\pandoc@box}%
  \ifdim\@tempb\p@<\@tempa\p@\let\@tempa\@tempb\fi% select the smaller of both
  \ifdim\@tempa\p@<\p@\scalebox{\@tempa}{\usebox\pandoc@box}%
  \else\usebox{\pandoc@box}%
  \fi%
}
% Set default figure placement to htbp
\def\fps@figure{htbp}
\makeatother
\setlength{\emergencystretch}{3em} % prevent overfull lines
\providecommand{\tightlist}{%
  \setlength{\itemsep}{0pt}\setlength{\parskip}{0pt}}
\usepackage{bookmark}
\IfFileExists{xurl.sty}{\usepackage{xurl}}{} % add URL line breaks if available
\urlstyle{same}
\hypersetup{
  pdftitle={Exercises 6 - solutions},
  pdfauthor={Maria Cuellar},
  hidelinks,
  pdfcreator={LaTeX via pandoc}}

\title{Exercises 6 - solutions}
\author{Maria Cuellar}
\date{2025-10-21}

\begin{document}
\maketitle

Note: use city\_crime\_spending.csv.

\begin{Shaded}
\begin{Highlighting}[]
\FunctionTok{library}\NormalTok{(tidyverse)}
\NormalTok{dat }\OtherTok{\textless{}{-}} \FunctionTok{read\_csv}\NormalTok{(}\StringTok{"/Users/mariacuellar/Github/crim\_data\_analysis/data/city\_crime\_spending.csv"}\NormalTok{)}
\end{Highlighting}
\end{Shaded}

\textbf{1. Question} For two quantitative variables (x and y): Write
down the requirements for each of these steps.

\begin{itemize}
\tightlist
\item
  Do visual EDA: Draw a scatterplot and describe the association, form,
  strengh, and outliers.
\item
  Do quantitative EDA: Check whether the two variables are quantitative,
  whether their relationship is straight enough, and whether there are
  any outliers.
\item
  Fit a linear model: Check assumptions of linear regression.
\item
  Test assumptions of linear regression: Linear relationship between x
  and y, Independence between observations, Homoscedasticity of errors,
  and Normality of y for a given x. Also, check that there are no
  influential outliers.
\item
  Interpret coefficient of x to do inference: Only if the assumptions of
  linear regression are satisfied.
\end{itemize}

\textbf{Instructions} Fit a simple linear regression for
police\_spending regressed onto population.

\textbf{2. Question} What is the null hypothesis for doing inference?
What is the research question you could ask here?

The null hypothesis is that there is no relationship between population
and police\_spending.

\textbf{3. Question} What is the coefficient of determination? (The one
called Multiple R-squared. The other one has a penalty for adding more
covariates, so we won't need it here.) What does this mean?

I fit a linear model:

\begin{Shaded}
\begin{Highlighting}[]
\NormalTok{out }\OtherTok{\textless{}{-}} \FunctionTok{lm}\NormalTok{(police\_spending }\SpecialCharTok{\textasciitilde{}}\NormalTok{ population, }\AttributeTok{data=}\NormalTok{dat)}
\end{Highlighting}
\end{Shaded}

And look at the R squared.

\begin{Shaded}
\begin{Highlighting}[]
\FunctionTok{summary}\NormalTok{(out)}
\end{Highlighting}
\end{Shaded}

\begin{verbatim}
## 
## Call:
## lm(formula = police_spending ~ population, data = dat)
## 
## Residuals:
##      Min       1Q   Median       3Q      Max 
## -148.345  -25.665   -2.046   20.840  168.405 
## 
## Coefficients:
##              Estimate Std. Error t value Pr(>|t|)    
## (Intercept) 9.078e+01  3.800e+00   23.89   <2e-16 ***
## population  3.186e-04  6.558e-06   48.59   <2e-16 ***
## ---
## Signif. codes:  0 '***' 0.001 '**' 0.01 '*' 0.05 '.' 0.1 ' ' 1
## 
## Residual standard error: 42.17 on 498 degrees of freedom
## Multiple R-squared:  0.8258, Adjusted R-squared:  0.8254 
## F-statistic:  2360 on 1 and 498 DF,  p-value: < 2.2e-16
\end{verbatim}

The (Multiple) R-squared is 0.8258. This means that there is a strong
association between population and police\_spending.

\textbf{4. Question} Test the assumptions of the linear regression.

First, draw a scatterplot.

\begin{Shaded}
\begin{Highlighting}[]
\NormalTok{dat }\SpecialCharTok{\%\textgreater{}\%} \FunctionTok{ggplot}\NormalTok{(}\FunctionTok{aes}\NormalTok{(}\AttributeTok{x=}\NormalTok{population, }\AttributeTok{y=}\NormalTok{police\_spending)) }\SpecialCharTok{+} \FunctionTok{geom\_point}\NormalTok{()}
\end{Highlighting}
\end{Shaded}

\pandocbounded{\includegraphics[keepaspectratio]{Exercises-6---solutions_files/Exercises-6---solutions_files/figure-latex/visual-EDA-1.pdf}}

For quantitative EDA, I check the conditions for correlation: x and y
are quantitative variables, the relationship between x and y is straight
enough, and there are no noticeable outliers. Therefore, I go ahead and
calculate the correlation.

\begin{Shaded}
\begin{Highlighting}[]
\NormalTok{dat }\SpecialCharTok{\%\textgreater{}\%} \FunctionTok{summarize}\NormalTok{(}\AttributeTok{correlation =} \FunctionTok{cor}\NormalTok{(population, police\_spending))}
\end{Highlighting}
\end{Shaded}

\begin{verbatim}
## # A tibble: 1 x 1
##   correlation
##         <dbl>
## 1       0.909
\end{verbatim}

I draw the diagnostic plots:

\begin{Shaded}
\begin{Highlighting}[]
\FunctionTok{par}\NormalTok{(}\AttributeTok{mfrow=}\FunctionTok{c}\NormalTok{(}\DecValTok{2}\NormalTok{,}\DecValTok{2}\NormalTok{))}
\FunctionTok{plot}\NormalTok{(out)}
\end{Highlighting}
\end{Shaded}

\pandocbounded{\includegraphics[keepaspectratio]{Exercises-6---solutions_files/Exercises-6---solutions_files/figure-latex/diagnostic-plots-1.pdf}}

\begin{Shaded}
\begin{Highlighting}[]
\FunctionTok{par}\NormalTok{(}\AttributeTok{mfrow=}\FunctionTok{c}\NormalTok{(}\DecValTok{1}\NormalTok{,}\DecValTok{1}\NormalTok{))}
\end{Highlighting}
\end{Shaded}

And then test assumptions:

\begin{enumerate}
\def\labelenumi{\alph{enumi})}
\tightlist
\item
  Linear relationship between x and y: There is a linear relationship
  between x and y, shown in the scatterplot.
\item
  Independence between observations: I cannot test for independence
  between observations, but I don't see any evidence of clumping to
  suggest strong dependencies.
\item
  Homoscedasticity of errors: There is a clear difference in variance of
  the errors shown in the scale-location plot, as well as in the
  residuals vs.~fitted plot and the scatterplot. This assumption is NOT
  met.
\item
  Normality of y for a given x: The q-q plot shows the points mostly on
  the diagonal, except for the ends. This shows me the assumption is
  mostly (sufficiently) met.
\item
  There are no influential outliers.
\end{enumerate}

\begin{itemize}
\tightlist
\item
  For visual EDA: I draw a scatterplot. The scatterplot shows that the
  two variables have a positive association, a linear form, the
  relationship is strong (although I do notice changing variance
  throughout ``population''), and there are no clear outliers.
\end{itemize}

\textbf{5. Question} Should you interpret the coefficients? If so, go
ahead and interpret the coefficients. If not, then give a satatement for
why you should not interpret them.

Since the assumptions of linear regression are NOT fully met, I will not
interpret the coefficients.

\end{document}
